\chapter{Conclusão} 

\label{cap6_conclusao}

O presente trabalho teve como objetivo desenvolver um protótipo experimental de um sistema fotovoltaico automatizado de baixo custo para avaliar a influência do reposicionamento automático do painel fotovoltaico na geração de energia elétrica, bem como desenvolver e validar um modelo computacional no ambiente MATLAB/Simulink por meio da comparação entre resultados simulados e experimentais.

Os resultados experimentais demonstraram que o sistema automatizado apresentou desempenho superior ao sistema com painel fixo durante todo o período de monitoramento. A análise das curvas de potência e da energia acumulada evidenciou que o reposicionamento automático do painel proporcionou um ganho energético médio de aproximadamente 29,69\%, confirmando que a estratégia de orientação adotada aumenta o aproveitamento da radiação solar incidente quando comparada à utilização de um painel mantido em posição fixa.

Na simulação computacional, o modelo desenvolvido no MATLAB/Simulink indicou um ganho energético de aproximadamente 23,94\% para o sistema automatizado em relação ao sistema com painel fixo. Embora esse valor seja inferior ao obtido experimentalmente, ambos apresentam a mesma tendência de aumento da geração de energia proporcionada pelo reposicionamento automático do painel. A proximidade entre os resultados simulados e experimentais demonstra que o modelo foi capaz de representar satisfatoriamente o comportamento geral do sistema, atendendo ao objetivo de validação proposto neste trabalho.

O protótipo desenvolvido cumpriu sua finalidade como plataforma experimental para aquisição de dados e validação do modelo computacional. A integração entre Arduino Uno, sensores de tensão e corrente, servomotor, módulo RTC, display LCD e cartão microSD possibilitou a implementação de um sistema capaz de realizar automaticamente o monitoramento das variáveis elétricas, o controle do posicionamento do painel fotovoltaico e o armazenamento dos dados experimentais.

Entretanto, os resultados devem ser analisados considerando as limitações inerentes à pesquisa. O protótipo foi desenvolvido exclusivamente para fins experimentais, utilizando um módulo fotovoltaico de pequena potência e uma estrutura mecânica simplificada, não representando uma solução destinada à aplicação em sistemas fotovoltaicos comerciais. Além disso, não foram avaliados aspectos relacionados à durabilidade mecânica da estrutura, aos custos de implantação e manutenção, ao consumo energético dos componentes eletrônicos utilizados na automação e à viabilidade técnico-econômica da solução proposta. Os ganhos energéticos apresentados referem-se exclusivamente ao aumento da energia produzida pelo painel fotovoltaico em decorrência do reposicionamento automático, não considerando o consumo dos dispositivos responsáveis pela automação do sistema. Também não foram incorporadas ao modelo computacional medições reais de irradiância e temperatura durante os ensaios experimentais, uma vez que o objetivo do estudo concentrou-se na análise do ganho de energia proporcionado pelo reposicionamento automático do painel.

Apesar dessas limitações, os resultados obtidos fornecem evidências experimentais de que a orientação automática do painel fotovoltaico contribui para o aumento da energia gerada ao longo do dia. Além disso, demonstram que o modelo computacional desenvolvido representa adequadamente o comportamento observado no protótipo físico, podendo ser utilizado como ferramenta de apoio em análises e estudos comparativos de sistemas fotovoltaicos automatizados.

Conclui-se, portanto, que os objetivos propostos foram plenamente alcançados. Foi possível desenvolver o protótipo experimental, implementar o sistema de aquisição e controle, construir e validar o modelo computacional e comparar os resultados simulados com os dados obtidos em condições reais de operação. Os resultados evidenciam que a utilização de sistemas embarcados associados a técnicas de automação constitui uma alternativa tecnicamente viável para aumentar o aproveitamento da energia solar em aplicações experimentais. Assim, este trabalho contribui para o desenvolvimento de pesquisas relacionadas a sistemas fotovoltaicos automatizados, modelagem computacional, geração distribuída e energias renováveis, fornecendo uma base experimental que poderá ser utilizada em estudos futuros voltados ao aperfeiçoamento de sistemas de rastreamento solar e à avaliação de sua viabilidade em aplicações de maior escala.